\chapter{Class templates in \texttt{Dwm::What} namespace}
There are two class templates of interest in the {\em dwmwhat} package:
\texttt{Dwm::What::SegmentedLiteral} and \texttt{Dwm::What::Info}.
We saw simple use of the \texttt{Dwm::What::Info} class template in
the introduction, and it is likely the only class template most users
will use.

However, \texttt{Dwm::What::Info} inherits from
\texttt{Dwm::What::SegmentedLiteral}, and I'll introduce them in
base to derived order.  
\texttt{Dwm::What::SegmentedLiteral} has no constraints on number
of segments or choice of delimiter (an empty delimiter is permissible),
and all of the segments are user defined.

\section{\texttt{Dwm::What::SegmentedLiteral}}
\texttt{Dwm::What::SegmentedLiteral} builds a string literal from
a given delimiter (placed between segments) and one or more segments
(given as string literals).  All arguments to the constructor are
string literals, and used to deduce the template parameters.

An empty delimiter is allowed, as are empty segments.

\subsection{Synopsis}

\begin{verbatim}
template <std::size_t DelimLen, std::size_t FirstLen, std::size_t ...SegLen>
class SegmentedLiteral
{
public:
   consteval SegmentedLiteral(const char (&delim)[DelimLen],
                              const char (&firstSeg)[FirstLen],
                              const char (&...moreSegs)[SegLen]);
   constexpr std::string_view view() const noexcept;
   constexpr std::size_t num_segments() const noexcept;      
   constexpr std::string_view nth(std::size_t n) const noexcept;
};
\end{verbatim}

\subsection{Member functions}

\begin{minipage}{\linewidth}
\subsubsection{Constructor}
The \texttt{Dwm::What::SegmentedLiteral} constructor requires a
delimiter (a string literal) and at least one segment.  All segments
are given as string literals.

\begin{description}[noitemsep,align=left,leftmargin=*,labelsep=10pt,itemindent=-\leftmargin,style=nextline]
  \item[\texttt{delim}] The delimiter that will be placed between segments,
  which may be empty.
  \item[\texttt{firstSeg}] The first segment.
  \item[\texttt{moreSegs}] The remaining segments.
\end{description}
\end{minipage}

\subsubsection{Accessors}
\begin{description}[noitemsep,align=left,leftmargin=*,labelsep=10pt,itemindent=-\leftmargin,style=nextline]
  \item[\texttt{view()}] returns a \texttt{std::string\_view} of the entire
  constructed string literal, sans the terminating null.
  \item[\texttt{num\_segments()}] returns the number of segments.
  \item[\texttt{nth(std::size\_t n)}] returns the \texttt{n}th segment (0 indexed).
\end{description}

\subsection{Examples}

\begin{minipage}{\linewidth}
\input{SegmentedLiteralEx01.tex}
\end{minipage}

\begin{verbatim}
Dwm::What::SegmentedLiteral("|", "@(#)", "DwmWhat", "1.0.0",
                            "Copyright Daniel McRobb 2025", "mcplex.net");
\end{verbatim}

\begin{center}
\noindent\begin{tikzpicture}[
    start chain = A going right,
    node distance = 0pt, % Ensure no gap between bit nodes
    byte node/.style = {draw=lightgray, minimum height=20pt, minimum width=(\textwidth)/56, text width=\dimexpr(\textwidth)/56\relax, text height=16pt, text depth=4pt, font=\ttfamily, anchor=south, align=center, inner sep=0pt, outer sep=0pt, on chain=A}
  ]
  \def\mybytestring{{@},{(},{\#},{)},{|},D,w,m,W,h,a,t,{|},1,.,0,.,0,{|},C,o,p,y,r,i,g,h,t,{ },D,a,n,i,e,l,{ },M,c,R,o,b,b,{ },2,0,2,5,{|},m,c,p,l,e,x,.,n,e,t}
  % Draw the byte nodes and indices
  \setcounter{lastbyte}{0}
  \foreach \byte [count=\i] in \mybytestring {
    \node[byte node, fill=white] (byte-\i) {\byte};
    \stepcounter{lastbyte}
  }
  \node[byte node, fill=lightgray!50] (termnull) {$\emptyset$};
  
  % and a label for the byte nodes
  \ifnum \value{lastbyte} > 0
    \draw [decorate, decoration={brace, amplitude=10pt}] (byte-1.north west) -- (byte-\thelastbyte.north east)
    node [above=8pt, midway, font=\scriptsize] {\texttt{view()}};
    \draw [decorate, decoration={brace, mirror, amplitude=10pt}] (byte-1.south west) -- (byte-4.south east)
    node [below=8pt, midway, font=\scriptsize] {\texttt{nth(0)}};
    \draw [decorate, decoration={brace, mirror, amplitude=10pt}] (byte-6.south west) -- (byte-12.south east)
    node [below=8pt, midway, font=\scriptsize] {\texttt{nth(1)}};
    \draw [decorate, decoration={brace, mirror, amplitude=10pt}] (byte-14.south west) -- (byte-18.south east)
    node [below=8pt, midway, font=\scriptsize] {\texttt{nth(2)}};
    \draw [decorate, decoration={brace, mirror, amplitude=10pt}] (byte-20.south west) -- (byte-47.south east)
    node [below=8pt, midway, font=\scriptsize] {\texttt{nth(3)}};
    \draw [decorate, decoration={brace, mirror, amplitude=10pt}] (byte-49.south west) -- (byte-58.south east)
    node [below=8pt, midway, font=\scriptsize] {\texttt{nth(4)}};

  \draw [decorate, decoration={brace, mirror, amplitude=10pt}] (byte-5.south west) -- (byte-5.south east) 
  node [below=8pt, midway, rotate=270, anchor=east, xshift=3em, yshift=.5pt, font=\scriptsize] {delim};
  \draw [decorate, decoration={brace, mirror, amplitude=10pt}] (byte-13.south west) -- (byte-13.south east) 
  node [below=8pt, midway, rotate=270, anchor=east, xshift=3em, yshift=.5pt, font=\scriptsize] {delim};
  \draw [decorate, decoration={brace, mirror, amplitude=10pt}] (byte-19.south west) -- (byte-19.south east) 
  node [below=8pt, midway, rotate=270, anchor=east, xshift=3em, yshift=.5pt, font=\scriptsize] {delim};
  \draw [decorate, decoration={brace, mirror, amplitude=10pt}] (byte-48.south west) -- (byte-48.south east) 
  node [below=8pt, midway, rotate=270, anchor=east, xshift=3em, yshift=.5pt, font=\scriptsize] {delim};

  \fi
\end{tikzpicture}
\captionof{figure}{\texttt{Dwm::What::SegmentedLiteral} layout}
\end{center}


\pagebreak

\section{\texttt{Dwm::What::Info}}
\texttt{Dwm::What::Info} is just a \texttt{Dwm::What::SegmentedLiteral}
with seven segments and a specific delimiter (not user-specified).
The first segment (\texttt{"@(\#)"}) is automatically added, and named
accessors are provided for the remaining six segments.

\begin{minipage}{\linewidth}
\subsection{Synopsis}

\begin{small}
\begin{verbatim}
template <std::size_t TypeLen, std::size_t StatusLen, std::size_t NameLen,
           std::size_t VersionLen, std::size_t CopyrightLen, std::size_t OtherLen>
class Info
   : public SegmentedLiteral<sizeof(DWM_WHAT_DELIM), sizeof("@(#)"), TypeLen,
                             StatusLen, NameLen, VersionLen, CopyrightLen, OtherLen>
{
public:
   consteval Info(const char (&type)[TypeLen], const char (&status)[StatusLen],
                  const char (&name)[NameLen], const char (&version)[VersionLen],
                  const char (&copyright)[CopyrightLen], const char (&other)[OtherLen]);

   constexpr std::string_view type() const noexcept;
   constexpr std::string_view status() const noexcept;
   constexpr std::string_view name() const noexcept;
   constexpr std::string_view version() const noexcept;
   constexpr std::string_view copyright() const noexcept;
   constexpr std::string_view other() const noexcept;
   constexpr std::string_view data_view() const noexcept;
   constexpr std::string_view view() const noexcept;
};
\end{verbatim}
\end{small}
\end{minipage}

\subsection{Constructor}
\subsubsection{constructor arguments}
\begin{description}[noitemsep,align=left,leftmargin=*,labelsep=10pt,itemindent=-\leftmargin,style=nextline]
\item[\texttt{type}] The type of the package.  By convention, one or more of
  \texttt{DWM\_WHAT\_TYPE\_EXE}, \texttt{DWM\_WHAT\_TYPE\_LIB},
  \texttt{DWM\_WHAT\_TYPE\_HDR} and \texttt{DWM\_WHAT\_TYPE\_DOC}.  These are
  just macros that expand to string literals.  See \autoref{table:Macros}.
\item[\texttt{status}] The status of the package.  By convention, one of
  \texttt{DWM\_WHAT\_STATUS\_REL}, \texttt{DWM\_WHAT\_STATUS\_RC} or
  \texttt{DWM\_WHAT\_STATUS\_DEV}.  These are just macros that expand to
  string literals.  See \autoref{table:Macros}.
  \item[\texttt{name}] The name of the package.
  \item[\texttt{version}] The version of the package.
  \item[\texttt{copyright}] The copyright.
  \item[\texttt{other}] Any other additional information.
\end{description}

\begin{minipage}{\linewidth}
\subsection{Accesssors}
\begin{description}[noitemsep,align=left,leftmargin=*,labelsep=10pt,itemindent=-\leftmargin,style=nextline]
\item[\texttt{type()}] Returns the type of the package.  This is the same that was
  provided as the first argument to the constructor.
\item[\texttt{status()}] Returns the status of the package.  This is the same that was
  provided as the second argument to the constructor.
\item[\texttt{name()}] Returns the name of the package.  This is the same that was
  provided as the third argument to the constructor.
\item[\texttt{version()}] Returns the version of the package.  This is the same that
  was provided as the fourth argument to the constructor.
\item[\texttt{copyright()}] Returns the copyright.  This is the same that was provided
  as the fifth argument to the constructor.
\item[\texttt{other()}] Returns any additional information that was provided.  This
  is the same as the sixth (final) argument to the constructor.
\end{description}
\end{minipage}

\subsection{Equivalence to \texttt{Dwm::What::SegmentedLiteral}}
\begin{minipage}{\linewidth}
The following will produce the equivalent string literal.  This makes
it obvious that \texttt{Dwm::What::Info} is just a
\texttt{Dwm::What::SegmentedLiteral} with a specific delimiter
(\texttt{"\textbackslash{xC2}\textbackslash{xA0}"}) and an
automatically added first segment (\texttt{"@(\#)"}).

\begin{Verbatim}
   SegmentedLiteral(\textcolor{blue}{"\textbackslash{xC2}\textbackslash{xA0}"}, \textcolor{blue}{"@(\#)"},
                    DWM_WHAT_TYPE_LIB, DWM_WHAT_STATUS_REL, "DwmWhat",
                    "1.0.0", "Daniel McRobb 2025", "mcplex.net");

               Info(DWM_WHAT_TYPE_LIB, DWM_WHAT_STATUS_REL, "DwmWhat",
                    "1.0.0", "Daniel McRobb 2025", "mcplex.net");
\end{Verbatim}
\end{minipage}

\subsection{Examples}
\begin{minipage}{\linewidth}
This call to the \texttt{Dwm::What::Info} constructor:
\begin{verbatim}
   Dwm::What::Info(DWM_WHAT_TYPE_EXE DWM_WHAT_TYPE_HDR, DWM_WHAT_STATUS_REL,
                   "DwmWhat", "1.0.0", "Daniel McRobb 2025", "mcplex.net");
\end{verbatim}

Would produce a 67 byte string literal at compile time with the layout shown
in \autoref{figure:InfoEx01}.  Note this figure shows what would be returned
by accessors inherited from the \texttt{Dwm::What::SegmentedLiteral} base
class (in black) and the accessors from the \texttt{Dwm::What::Info} class
(in \textcolor{blue}{blue}).

\begin{center}
  \begin{figure}[H]
\noindent\begin{tikzpicture}[
  start chain = A going right,
  node distance = 0pt, % Ensure no gap between bit nodes
  byte node/.style = {draw=lightgray, minimum height=20pt, text height=16pt, text depth=4pt, font=\ttfamily, anchor=south, align=center, inner sep=0pt, outer sep=0pt, on chain=A}
]
\newlength{\onebyte}
\setlength{\onebyte}{1.7ex}

\node[byte node,minimum width=4\onebyte,text width=4\onebyte,fill=white] (seg1) {@(\#)};
\node[byte node,minimum width=\onebyte,text width=\onebyte, fill=red!15,font=\scriptsize\ttfamily, text depth=0pt] (delim1a) {\rotatebox{90}{0xC2}};
\node[byte node,minimum width=\onebyte,text width=\onebyte, fill=red!15,font=\scriptsize\ttfamily, text depth=0pt] (delim1b) {\rotatebox{90}{0xA0} };
\node[byte node,minimum width=7\onebyte,text width=7\onebyte,fill=white] (seg2) {\emoji{robot}\#};
\node[byte node,minimum width=\onebyte,text width=\onebyte, fill=red!15,font=\scriptsize\ttfamily, text depth=0pt] (delim2a) {\rotatebox{90}{0xC2}};
\node[byte node,minimum width=\onebyte,text width=\onebyte, fill=red!15,font=\scriptsize\ttfamily, text depth=0pt] (delim2b) {\rotatebox{90}{0xA0} };
\node[byte node,minimum width=3\onebyte,text width=3\onebyte,fill=white] (seg3) {\emoji{white-check-mark}};
\node[byte node,minimum width=\onebyte,text width=\onebyte, fill=red!15,font=\scriptsize\ttfamily, text depth=0pt] (delim3a) {\rotatebox{90}{0xC2}};
\node[byte node,minimum width=\onebyte,text width=\onebyte, fill=red!15,font=\scriptsize\ttfamily, text depth=0pt] (delim3b) {\rotatebox{90}{0xA0} };
\node[byte node,minimum width=7\onebyte,text width=7\onebyte,fill=white] (seg4) {DwmWhat};
\node[byte node,minimum width=\onebyte,text width=\onebyte, fill=red!15,font=\scriptsize\ttfamily, text depth=0pt] (delim4a) {\rotatebox{90}{0xC2}};
\node[byte node,minimum width=\onebyte,text width=\onebyte, fill=red!15,font=\scriptsize\ttfamily, text depth=0pt] (delim4b) {\rotatebox{90}{0xA0} };
\node[byte node,minimum width=5\onebyte,text width=5\onebyte,fill=white] (seg5) {1.0.0};
\node[byte node,minimum width=\onebyte,text width=\onebyte, fill=red!15,font=\scriptsize\ttfamily, text depth=0pt] (delim5a) {\rotatebox{90}{0xC2}};
\node[byte node,minimum width=\onebyte,text width=\onebyte, fill=red!15,font=\scriptsize\ttfamily, text depth=0pt] (delim5b) {\rotatebox{90}{0xA0} };
\node[byte node,minimum width=18\onebyte,text width=18\onebyte,fill=white] (seg6) {Daniel McRobb 2025};
\node[byte node,minimum width=\onebyte,text width=\onebyte, fill=red!15,font=\scriptsize\ttfamily, text depth=0pt] (delim6a) {\rotatebox{90}{0xC2}};
\node[byte node,minimum width=\onebyte,text width=\onebyte, fill=red!15,font=\scriptsize\ttfamily, text depth=0pt] (delim6b) {\rotatebox{90}{0xA0} };
\node[byte node,minimum width=10\onebyte,text width=10\onebyte,fill=white] (seg7) {mcplex.net};
\node[byte node,minimum width=\onebyte,text width=\onebyte,fill=lightgray!50] (termnull) {$\emptyset$};

\node (nth0) [below=of seg1.south, yshift=-20pt,font=\scriptsize]
      {\texttt{nth(0)}};
\draw [decorate, decoration={brace, mirror, amplitude=10pt}]
  (seg2.south west) -- (seg2.south east)
  node [below=8pt, midway, font=\scriptsize, text=blue] {\texttt{type()}};
\node (nth1) [below=of seg2.south, yshift=-20pt,font=\scriptsize]
      {\texttt{nth(1)}};
\draw [decorate, decoration={brace, mirror, amplitude=10pt}]
  (seg3.south west) -- (seg3.south east)
  node [below=8pt, midway, font=\scriptsize, text=blue] {\texttt{status()}};
\node (nth2) [below=of seg3.south, yshift=-20pt,font=\scriptsize]
      {\texttt{nth(2)}};
\draw [decorate, decoration={brace, mirror, amplitude=10pt}]
  (seg4.south west) -- (seg4.south east)
  node [below=8pt, midway, font=\scriptsize, text=blue] {\texttt{name()}};
\node (nth3) [below=of seg4.south, yshift=-20pt,font=\scriptsize]
      {\texttt{nth(3)}};
\draw [decorate, decoration={brace, mirror, amplitude=10pt}]
  (seg5.south west) -- (seg5.south east)
  node [below=8pt, midway, font=\scriptsize, text=blue] {\texttt{version()}};
\node (nth4) [below=of seg5.south, yshift=-20pt,font=\scriptsize]
      {\texttt{nth(4)}};
\draw [decorate, decoration={brace, mirror, amplitude=10pt}]
  (seg6.south west) -- (seg6.south east)
  node [below=8pt, midway, font=\scriptsize, text=blue] {\texttt{copyright()}};
\node (nth6) [below=of seg6.south, yshift=-20pt,font=\scriptsize]
      {\texttt{nth(5)}};
\draw [decorate, decoration={brace, mirror, amplitude=10pt}]
  (seg7.south west) -- (seg7.south east)
  node [below=8pt, midway, font=\scriptsize, text=blue] {other()};
\node (nth8) [below=of seg7.south, yshift=-20pt,font=\scriptsize]
      {\texttt{nth(6)}};
      
\draw [decorate, decoration={brace, amplitude=10pt}]
  (seg2.north west) -- (seg7.north east)
  node [above=8pt, midway, font=\scriptsize, text=blue] {\texttt{data\_view()}};

\draw [decorate, decoration={brace, amplitude=40pt}]
  (seg1.north west) -- (seg7.north east)
  node [above=40pt, midway, font=\scriptsize] {\texttt{view()}};

\end{tikzpicture}
\caption{\texttt{Dwm::What::Info} layout}
\label{figure:InfoEx01}
\end{figure}
\end{center}
\end{minipage}

The 67 bytes:
\begin{tabular}{ll}
  @(\#)            &  4 bytes \\
  delim            &  2 bytes \\
  type             &  7 bytes \\
  delim            &  2 bytes \\
  status           &  3 bytes \\
  delim            &  2 bytes \\
  name             &  7 bytes \\
  delim            &  2 bytes \\
  version          &  5 bytes \\
  delim            &  2 bytes \\
  copyright        & 18 bytes \\
  delim            &  2 bytes \\
  other            & 10 bytes \\
  null             &  1 byte \\ \hline
  \textbf{Total} & \textbf{67 bytes}
\end{tabular}

And {\em dwmwhat} would print this when scanning a compiled program
containing such an instance of \texttt{Dwm::What::Info}:

\texttt{
\emoji{robot}\# \emoji{white-check-mark} DwmWhat 1.0.0 Daniel McRobb 2025 mcplex.net
}


\subsection{\texttt{SegmentedLiteral} and \texttt{Info} equivalents}

\begin{minipage}{\linewidth}
  All of the following will produce the same string, as show in
  \autoref{figure:SLInfoEquiv01}.  The first two are using macros from
  \texttt{DwmWhatInfo.hh} (\autoref{table:Macros}) in constructor
  arguments.  The next two are repeats of the first two, with macro
  substitution.

\begin{small}
\begin{verbatim}
Dwm::What::SegmentedLiteral(DWM_WHAT_SYM_NB_SPACE, "@(#)", DWM_WHAT_TYPE_LIB,
                            DWM_WHAT_STATUS_REL, "dwmwhat", "1.0.0",
                            "Daniel McRobb", "mcplex.net");
\end{verbatim}
\color{blue}
\begin{verbatim}
            Dwm::What::Info(DWM_WHAT_TYPE_LIB, DWM_WHAT_STATUS_REL, "dwmwhat",
                            "1.0.0", "Daniel McRobb", "mcplex.net");
\end{verbatim}
\color{black}

\begin{verbatim}
Dwm::What::SegmentedLiteral("\xC2\xA0", "@(#)", "\xF0\x9F\x93\x9A", "\xE2\x9C\x85",
                            "dwmwhat", "1.0.0", "Daniel McRobb", "mcplex.net");
\end{verbatim}
\color{blue}
\begin{verbatim}
            Dwm::What::Info("\xF0\x9F\x93\x9A", "\xE2\x9C\x85", "dwmwhat",
                            "1.0.0", "Daniel McRobb", "mcplex.net");
\end{verbatim}
\color{black}
\end{small}

\begin{verbatim}
Dwm::What::SegmentedLiteral("\xC2\xA0", "@(#)", DWM_WHAT_TYPE_LIB,
                            DWM_WHAT_STATUS_REL, "DwmWhat", "1.0.0",
                            "Daniel McRobb 2025", "mcplex.net");
\end{verbatim}

\begin{center}
\noindent\begin{tikzpicture}[
    start chain = A going right,
    node distance = 0pt, % Ensure no gap between bit nodes
    byte node/.style = {draw=lightgray, minimum height=20pt, minimum width=(\textwidth)/56, text width=\dimexpr(\textwidth)/56\relax, text height=16pt, text depth=4pt, font=\ttfamily, anchor=south, align=center, inner sep=0pt, outer sep=0pt, on chain=A}
]
  \def\mypfxstring{{@},{(},{\#},{)}}
  \def\mydelimstring{0xC2,0xA0}
  \def\mysegonestring{0xF0,0x9F,0x93,0x9A}
  \def\mysegtwostring{0xE2,0x9C,0x85}
  \def\mysegthreestring{D,w,m,W,h,a,t}
  \def\mysegfourstring{1,.,0,.,0}
  \def\mysegfivestring{D,a,n,i,e,l,{ },M,c,R,o,b,b}
  \def\mysegsixstring{m,c,p,l,e,x,.,n,e,t}

  \foreach \byte [count=\i] in \mypfxstring {
    \node[byte node, fill=white] (byte-\i) {\byte};
    \stepcounter{lastbyte}
  }
  \draw [decorate, decoration={brace, mirror, amplitude=10pt}]
    (byte-1.south west) -- (byte-4.south east)
    node [below=8pt, midway, font=\scriptsize] {\texttt{nth(0)}};
  \draw [decorate, decoration={brace, amplitude=10pt}]
    (byte-1.north west) -- (byte-4.north east)
    node [above=8pt, midway] {@(\#)};
  \foreach \byte [count=\i] in \mydelimstring {
    \node[byte node, fill=red!15, font=\scriptsize\ttfamily, text depth=0pt] (byte-\i) {\rotatebox{90}{\byte}};
  }

  \foreach \byte [count=\i] in \mysegonestring {
    \node[byte node, fill=white, font=\scriptsize\ttfamily, text depth=0pt] (byte-\i) {\rotatebox{90}{\byte}};
    \stepcounter{lastbyte}
  }
  \draw [decorate, decoration={brace, mirror, amplitude=10pt}]
  (byte-7.south west) -- (byte-10.south east)
  node [below=8pt, midway, font=\scriptsize] {\texttt{nth(1)}};
  \draw (byte-7.south west) -- (byte-10.south east)
    node [below=20pt, midway, font=\scriptsize, text=blue] {\texttt{type()}};
  \draw [decorate, decoration={brace, amplitude=10pt}]
  (byte-7.north west) -- (byte-10.north east)
  node [above=8pt, midway] {\emoji{books}};
  \foreach \byte [count=\i] in \mydelimstring {
    \node[byte node, fill=red!15, font=\scriptsize\ttfamily, text depth=0pt] (byte-\i) {\rotatebox{90}{\byte}};
  }

  \foreach \byte [count=\i] in \mysegtwostring {
    \node[byte node, fill=white, font=\scriptsize\ttfamily, text depth=0pt] (byte-\i) {\rotatebox{90}{\byte}};
    \stepcounter{lastbyte}
  }
  \draw [decorate, decoration={brace, mirror, amplitude=10pt}]
  (byte-13.south west) -- (byte-15.south east)
  node [below=8pt, midway, font=\scriptsize] {\texttt{nth(2)}};
  \draw (byte-13.south west) -- (byte-15.south east)
    node [below=20pt, midway, font=\scriptsize, text=blue] {\texttt{status()}};
  \draw [decorate, decoration={brace, amplitude=10pt}]
  (byte-13.north west) -- (byte-15.north east)
  node [above=8pt, midway] {\emoji{white-check-mark}};
  \foreach \byte [count=\i] in \mydelimstring {
    \node[byte node, fill=red!15, font=\scriptsize\ttfamily, text depth=0pt] (byte-\i) {\rotatebox{90}{\byte}};
  }

  \foreach \byte [count=\i] in \mysegthreestring {
    \node[byte node, fill=white] (byte-\i) {\byte};
    \stepcounter{lastbyte}
  }
  \draw [decorate, decoration={brace, mirror, amplitude=10pt}]
    (byte-18.south west) -- (byte-24.south east)
    node [below=8pt, midway, font=\scriptsize] {\texttt{nth(3)}};
  \draw (byte-18.south west) -- (byte-24.south east)
    node [below=20pt, midway, font=\scriptsize, text=blue] {\texttt{name()}};
  \draw [decorate, decoration={brace, amplitude=10pt}]
    (byte-18.north west) -- (byte-24.north east)
    node [above=8pt, midway] {DwmWhat};
  \foreach \byte [count=\i] in \mydelimstring {
    \node[byte node, fill=red!15, font=\scriptsize\ttfamily, text depth=0pt] (byte-\i) {\rotatebox{90}{\byte}};
  }

  \foreach \byte [count=\i] in \mysegfourstring {
    \node[byte node, fill=white] (byte-\i) {\byte};
    \stepcounter{lastbyte}
  }
  \draw [decorate, decoration={brace, mirror, amplitude=10pt}]
    (byte-27.south west) -- (byte-31.south east)
    node [below=8pt, midway, font=\scriptsize] {\texttt{nth(4)}};
  \draw (byte-27.south west) -- (byte-31.south east)
    node [below=20pt, midway, font=\scriptsize, text=blue] {\texttt{version()}};
  \draw [decorate, decoration={brace, amplitude=10pt}]
    (byte-27.north west) -- (byte-31.north east)
    node [above=8pt, midway] {1.0.0};
  \foreach \byte [count=\i] in \mydelimstring {
    \node[byte node, fill=red!15, font=\scriptsize\ttfamily, text depth=0pt] (byte-\i) {\rotatebox{90}{\byte}};
  }

  \foreach \byte [count=\i] in \mysegfivestring {
    \node[byte node, fill=white] (byte-\i) {\byte};
    \stepcounter{lastbyte}
  }
  \draw [decorate, decoration={brace, mirror, amplitude=10pt}]
    (byte-34.south west) -- (byte-46.south east)
    node [below=8pt, midway, font=\scriptsize] {\texttt{nth(5)}};
  \draw (byte-34.south west) -- (byte-46.south east)
    node [below=20pt, midway, font=\scriptsize, text=blue] {\texttt{copyright()}};
  \draw [decorate, decoration={brace, amplitude=10pt}]
    (byte-34.north west) -- (byte-46.north east)
    node [above=8pt, midway] {Daniel McRobb};
  \foreach \byte [count=\i] in \mydelimstring {
    \node[byte node, fill=red!15, font=\scriptsize\ttfamily, text depth=0pt] (byte-\i) {\rotatebox{90}{\byte}};
  }

  \foreach \byte [count=\i] in \mysegsixstring {
    \node[byte node, fill=white] (byte-\i) {\byte};
    \stepcounter{lastbyte}
  }
  \draw [decorate, decoration={brace, mirror, amplitude=10pt}]
    (byte-49.south west) -- (byte-58.south east)
    node [below=8pt, midway, font=\scriptsize] {\texttt{nth(6)}};
  \draw (byte-49.south west) -- (byte-58.south east)
    node [below=20pt, midway, font=\scriptsize, text=blue] {\texttt{other()}};
  \draw [decorate, decoration={brace, amplitude=10pt}]
    (byte-49.north west) -- (byte-58.north east)
    node [above=8pt, midway] {mcplex.net};

  \node[byte node, fill=lightgray!50] (termnull) {$\emptyset$};
    
\end{tikzpicture}
\captionof{figure}{\texttt{Dwm::What::SegmentedLiteral} layout}
\end{center}


\end{minipage}
